\section{Scope, Chronology, and Evidence}

\subsection{Identity, names, and range}

Subject: Origen of Alexandria, called Adamantius (``man of adamant'':
Eusebius, \work{Church History} 6.14.10; Jerome, \work{Letter} 33; Photius,
\work{Bibliotheca} 118). Whole-life range used in this study: born c.~185
(inferred from ``not quite seventeen'' at his father's death in the
persecution of 202/203 and ``sixty-nine years of age'' at death under
Gallus; the two inferences do not perfectly cohere, allowing c.~184--186);
died c.~253/254, probably at Tyre. Principal places: Alexandria (to
231/232); Caesarea Maritima in Palestine (from 231/232); with attested
stays or visits at Rome, Arabia (including Bostra), Antioch, Athens and
Greece, Cappadocia, and Tyre. No civil record,
inscription, or contemporary archive attests any of it; the documentary
base is literary.

\subsection{Reader, question, and thesis boundary}

Written for a serious general reader, student, or researcher who wants one
coherent, evidence-controlled life. The governing questions are stated in
the opening section. The organizing synthesis---that Origen's achievement
and his condemnation must be held together, and that the question ``was
Origen anathematized?'' admits only a divided answer---is historical, not
ecclesiastical. This study does not adjudicate the self-mutilation report,
the canonical validity of Demetrius's proceedings or of the Caesarean
ordination, the doctrinal merits of the condemned propositions, the
authority of the 543 or 553 anathemas, or Origen's sanctity or eternal
state; it claims no theological or ecclesiastical review or approval.

\subsection{Corpus and citation conventions}

Origen's works are cited by conventional English titles with the book,
chapter, and section numbering of the Ante-Nicene Fathers translation,
vol.~4 (Buffalo, 1885) for \work{On First Principles} (with Rufinus's
preface), \work{Against Celsus}, and the \work{Letter to Gregory};
quotations of Eusebius, Jerome's \work{On Illustrious Men} and letters,
and Vincent of L\'erins follow the Nicene and Post-Nicene Fathers, series 2;
Gregory Thaumaturgus follows Ante-Nicene Fathers vol.~6; Photius follows
J.~H. Freese's 1920 translation; the conciliar texts follow H.~Percival's
NPNF2 vol.~14 with N.~Tanner's edition for the editorial status of the
fifteen anathemas. All quotations were checked against the online
transcriptions identified in the References on 2026-07-24. \work{On First
Principles} is quoted throughout in and as the Rufinian Latin layer (so
marked in the text); no Greek retroversion is asserted. Works not
independently inspected in this research (notably \work{Commentary on
Matthew} 15 and Epiphanius, \work{Panarion} 64) are cited only through the
named scholarship that reports them.

\subsection{Evidence classes}

\begin{evidencekey}
\evidenceclass{A}{Origen's own surviving works and letters: contemporary,
but exegesis, apologetic, and advocacy, not memoir; several survive only in
announcedly corrected Latin translation, and each claim is limited by its
transmission layer.}
\evidenceclass{N}{Near-contemporary external witness: Gregory Thaumaturgus's
farewell panegyric (praise by a grateful pupil); Porphyry's hostile notice
(preserved only by Eusebius).}
\evidenceclass{E}{Early received tradition: Eusebius (the biographical
spine, written c.~70--80 years after Origen's death by a partisan with
documentary access), Pamphilus's \work{Defense} (via Photius), Jerome,
Rufinus, Vincent of L\'erins.}
\evidenceclass{L}{Later tradition: Photius's ninth-century epitomes, the
Tyre tomb tradition, medieval and early-modern portraits.}
\evidenceclass{M}{Official ecclesial acts, limited to each act's actual
object: the 543 synodal anathemas, the 553 conciliar documents, Catechism
citations, Benedict XVI's 2007 audiences.}
\evidenceclass{K}{Named modern critical reconstruction (Diekamp, Prat,
Crouzel, de Lubac, McGuckin, Edwards, Price, Perrone and others), used to
its stated premises.}
\evidenceclass{S}{Source-grounded project synthesis, never attributed to a
source.}
\end{evidencekey}

\subsection{Chronology}

Dates below are claims with stated bases; regnal synchronisms are
Eusebius's unless noted. ``Convention'' marks customary dates this study
does not assert as established.

\begin{lifetimeline}
\lifeevent{c.~184--186}{Birth, probably at Alexandria}{Inference from HE
6.2.12 with 6.2.2, and HE 7.1 (E)}{Range only}
\lifeevent{202/203}{Persecution at Alexandria; Leonides beheaded; property
confiscated}{``Tenth year of Severus'' (HE 6.2.2; 6.1) (E)}{Secure}
\lifeevent{c.~203}{Origen, ``in his eighteenth year,'' takes over
catechetical instruction; Demetrius confirms}{HE 6.3 (E)}{Secure as
sequence; age from Eusebius}
\lifeevent{c.~203--210s}{Ascetic regime; pupil martyrs; sale of library;
self-mutilation report placed here}{HE 6.3--8 (E)}{Report contested; see
tradition audit}
\lifeevent{198--217}{Visit to Rome under Zephyrinus}{HE 6.14.10 (E)}{Secure
event; date range only}
\lifeevent{c.~215/216}{Flight to Caesarea amid ``a considerable war'';
lay preaching; recall by Demetrius}{HE 6.19.16--19 (E); Caracalla link
(K)}{Plausible}
\lifeevent{222--235}{Audience with Julia Mamaea at Antioch}{HE 6.21 (E);
within Alexander Severus's reign}{Secure event; date range only}
\lifeevent{before 231/232}{\work{On First Principles}, early commentaries,
at Alexandria}{HE 6.24 (E)}{Secure as relative order}
\lifeevent{c.~231}{Ordination at Caesarea en route to Greece}{HE 6.23.4
(E)}{Secure}
\lifeevent{231/232}{Alexandrian proceedings (banishment; deposition);
final removal to Caesarea; Heraclas succeeds; Demetrius dies}{HE 6.26 (E);
Photius, cod.~118 (E/L)}{Proceedings secure; details unresolved}
\lifeevent{c.~233--238}{Gregory (Thaumaturgus) studies at Caesarea; the
\work{Panegyric} on his departure}{Panegyric (N); HE 6.30 (E); five-year
span is later report (K)}{Plausible}
\lifeevent{235--238}{Maximin's persecution; \work{Exhortation to
Martyrdom} for Ambrose and Protoctetus}{HE 6.28 (E)}{Secure}
\lifeevent{c.~238--244}{Beryllus of Bostra recalled from error; Arabian
soul dispute settled}{HE 6.33; 6.37 (E)}{Secure as events; dates
approximate}
\lifeevent{c.~246--249}{Discourses taken by stenographers; \work{Against
Celsus}; letters to Philip and Severa}{HE 6.36 (``over sixty'') (E)}
{Plausible--secure}
\lifeevent{250--251}{Decian persecution: imprisonment, iron collar, stocks,
threats of fire; judge avoids killing him}{HE 6.39.5, from Origen's letters
(E/A)}{Secure}
\lifeevent{c.~253/254}{Death ``being sixty-nine years of age,'' under
Gallus; at Tyre, and buried there}{HE 7.1 (E); Jerome, \work{On
Illustrious Men} 54 (E); Pamphilus's martyrdom-at-Caesarea variant in
Photius 118 (E/L)}{Year approximate; place from Jerome; variant
preserved}
\lifeevent{c.~307--310}{Pamphilus and Eusebius, \work{Defense of Origen}}
{Photius 118 (L); Jerome 75 (E)}{Secure}
\lifeevent{c.~374--377}{Epiphanius's dossier against Origen
(\work{Panarion} 64)}{Named scholarship (K)}{Secure}
\lifeevent{393--404}{First Origenist controversy: Rufinus's translation
(398), Jerome's letters, Theophilus's synod (400), Anastasius}{Rufinus's
preface (E); Jerome, \work{Letter} 84 (E)}{Secure}
\lifeevent{543}{Justinian's edict and anathemas adopted by the standing
synod; patriarchal and papal assent}{Percival's text (M); Prat (K)}{Secure}
\lifeevent{553}{Constantinople II: anathema 11 lists Origen; fifteen
anathemas subscribed around, not in, the council}{Conciliar text (M);
Tanner note, Diekamp via Prat, Price (K)}{Facts secure; conciliar status
disputed}
\lifeevent{1941, 2012}{Tura papyri; Munich Psalm homilies identified}
{Named scholarship, BSB Cod.\ graec.\ 314 (K)}{Secure}
\lifeevent{2007}{Benedict XVI's two catecheses on Origen}{Vatican texts
(M)}{Secure}
\end{lifetimeline}

\subsection{Method, limits, and negative results}

Earliest controlled evidence governs; witnesses separated by centuries are
never silently merged; where they conflict (the mutilation report, the
death accounts, the status of the 553 anathemas), the conflict is stated
and left open. Geographic and linguistic limits: the study reads English
translations of Greek and Latin sources; no manuscript, papyrological,
epigraphic, or archaeological research was undertaken; \work{On First
Principles} is treated throughout as Rufinian Latin. Bounded negative
results, from searches conducted 2026-07-24: no feast, cult, canonization,
or liturgical commemoration of Origen was located in any period or rite;
no ancient account of his death by an eyewitness, and no surviving
identifiable tomb, was located; no acts of the Alexandrian proceedings
were located (none are extant). These are bounded, correctable search
results, not proofs of absence. Mutable online sources (Vatican website,
New Advent, CCEL, tertullian.org, papalencyclicals.net, the Stanford
Encyclopedia) were used as checked on 2026-07-24.

\subsection{Rights and review state}

Project prose is original; quotations of ancient authors follow the
public-domain Ante-Nicene Fathers (1885--1886) and Nicene and Post-Nicene
Fathers translations, Freese's 1920 Photius, and Percival's 1900 conciliar
volume, as identified in the References; the 1911 Catholic Encyclopedia is
public domain. Copyrighted modern texts (Tanner's conciliar edition, the
Stanford Encyclopedia entry, modern scholarship) are paraphrased or quoted
only in brief attributed snippets within ordinary scholarly quotation;
Vatican texts (Catechism, papal audiences) are paraphrased with citation.
This publication has received internal editorial checking only: no
independent specialist, theological, or ecclesiastical review, and no
rights opinion, is claimed.
